\chapter{Use of Generative AI}

\section{Brainstorming}
Since generative AI tools like ChatGPT can quickly provide answers, it was used to ask trends that are happening in the field of Blockchain. It was asked to provide recent published papers that were checked by the proponents. It would give several research titles with the authors and links. The proponents would look at each of them to ensure their legitimacy. Extraction of the content of the provided researches were not used as there was no guarantee that the AI would give correct results. Tasks like those are prone to hallucinations.  

Another usage was the creation of keywords that can be used to search for papers from the online databases found in the internet. This made the curation of papers easier and faster. This, however, did not replace the manual reading of the existing papers especially their RRL section to find keywords, know more about the topic, and look for possible mentioned references. 

\section{Research Proper}
During the writing of this paper, AI was used to automatically convert comma separated values into tables using latex syntax. This was done for convenience since the actual data was already created by the proponents. 

AI was also used in debugging error messages given by latex compiler during the compilation of the PDF format.